% ============================================================================
% Sentiment-Driven Market Microstructure: An Agent-Based Modeling Framework
% ============================================================================
% Author: Murad Farzulla
% Organization: Farzulla Research
% Version: 1.1.0
% Date: December 2025
% ============================================================================

\documentclass[11pt,twocolumn]{article}

\usepackage[a4paper, margin=0.9in]{geometry}
\usepackage{setspace}
\usepackage[T1]{fontenc}
\usepackage{lmodern}
\usepackage{graphicx}
\usepackage{float}
\usepackage{booktabs}
\usepackage{array}
\usepackage{amsmath}
\usepackage{amssymb}
% Algorithm packages removed - not needed for methodology paper
\usepackage[round,authoryear]{natbib}
\bibliographystyle{plainnat}
\usepackage{xcolor}
\definecolor{farzullaburgundy}{RGB}{128,0,32}
\definecolor{orcidgreen}{RGB}{166,206,57}
\usepackage{tcolorbox}
\usepackage{tikz}
\usepackage{scalerel}
\usepackage{enumitem}
\usepackage{url}
\usepackage[colorlinks=true,
            linkcolor=farzullaburgundy,
            citecolor=farzullaburgundy,
            urlcolor=farzullaburgundy,
            breaklinks=true,
            pdftitle={Sentiment-Driven Market Microstructure},
            pdfauthor={Murad Farzulla},
            pdfkeywords={agent-based model, market microstructure, sentiment analysis, cryptocurrency, Monte Carlo dropout}]{hyperref}

\usepackage{titlesec}
\titlespacing*{\section}{0pt}{1.5ex plus 0.5ex minus 0.2ex}{1ex plus 0.2ex}
\titlespacing*{\subsection}{0pt}{1.2ex plus 0.4ex minus 0.2ex}{0.8ex plus 0.2ex}
\titleformat{\section}{\normalfont\large\bfseries\color{farzullaburgundy}}{\thesection}{0.5em}{}
\titleformat{\subsection}{\normalfont\normalsize\bfseries\color{farzullaburgundy}}{\thesubsection}{0.5em}{}

\usepackage{fancyhdr}

\newcommand{\orcidicon}{\scalerel*{
    \begin{tikzpicture}[x=3ex,y=3ex]
    \draw[fill=orcidgreen] (0,0) circle (0.5);
    \end{tikzpicture}
}{\textrm{I}}}

\newcommand{\paperver}{1.1.0}
\newcommand{\paperdate}{December 2025}
\newcommand{\paperdoi}{10.5281/zenodo.17957746}

\newcommand{\metadatabox}[1]{%
\begin{tcolorbox}[
    colback=gray!5,
    colframe=farzullaburgundy,
    title=Publication Metadata,
    fonttitle=\bfseries,
    coltitle=white,
    colbacktitle=farzullaburgundy,
    width=\columnwidth,
    arc=2mm,
    boxrule=0.8pt
]
\small
\textbf{DOI:} \href{https://doi.org/#1}{\texttt{#1}}\\
\textbf{Version:} \paperver\\
\textbf{Date:} \paperdate\\
\textbf{License:} CC-BY-4.0
\end{tcolorbox}
}

\pagestyle{fancy}
\fancyhf{}
\fancyhead[L]{\small\href{https://farzulla.org}{farzulla.org}}
\fancyhead[R]{\small\href{https://doi.org/\paperdoi}{DOI: \paperdoi}}
\fancyfoot[C]{\small\thepage}
\fancyfoot[L]{\small Murad Farzulla}
\fancyfoot[R]{\small v\paperver~| \paperdate}
\renewcommand{\headrulewidth}{0.4pt}
\renewcommand{\footrulewidth}{0.4pt}

\fancypagestyle{firstpage}{
  \fancyhf{}
  \fancyfoot[C]{\small\thepage}
  \fancyfoot[L]{\small Murad Farzulla}
  \fancyfoot[R]{\small v\paperver~| \paperdate}
  \renewcommand{\headrulewidth}{0pt}
  \renewcommand{\footrulewidth}{0.4pt}
}

\begin{document}

\onecolumn
\setstretch{1.2}
\thispagestyle{firstpage}

\begin{center}
{\small\color{farzullaburgundy}\textbf{WORKING PAPER v\paperver} | \color{gray}Not peer-reviewed}
\end{center}
\vspace{0.5em}

\begin{center}
{\Large\bfseries Sentiment-Driven Market Microstructure}\\[0.5em]
{\large\itshape An Agent-Based Modeling Framework for Cryptocurrency Markets}\\[1em]
{\bfseries Murad Farzulla}\textsuperscript{1} \href{https://orcid.org/0009-0002-7164-8704}{\orcidicon\ \texttt{0009-0002-7164-8704}}\\[0.5em]
{\small\itshape \textsuperscript{1}\href{https://farzulla.org}{Farzulla Research}}\\[0.3em]
{\small \paperdate}\\[0.5em]
{\small Correspondence: \href{mailto:murad@farzulla.org}{murad@farzulla.org}}
\end{center}

\begin{abstract}
\noindent This paper presents a methodological framework for agent-based modeling (ABM) of cryptocurrency market microstructure with real-time sentiment integration, along with preliminary simulation results. The proposed architecture combines three novel components: (1) a Monte Carlo Dropout sentiment analyzer providing uncertainty-aware sentiment scores from social media streams, (2) a heterogeneous agent population including market makers, informed traders, noise traders, and arbitrageurs with sentiment-responsive behavior rules, and (3) a dynamic factor model (DFM) for regime detection enabling adaptive agent parameterization.

The framework is designed for deployment on distributed computing infrastructure (Kubernetes/K3s) with Kafka-based event streaming for real-time data ingestion from cryptocurrency exchanges (Binance) and social platforms (Reddit). Sentiment processing uses CryptoBERT---fine-tuned on 3.2 million cryptocurrency social media posts---with Monte Carlo Dropout to quantify both aleatoric uncertainty (irreducible noise in sentiment signals) and epistemic uncertainty (model confidence), enabling agents to weight sentiment information appropriately.

We detail the technical architecture, agent behavior specifications, and calibration methodology for matching simulated stylized facts to empirical observations. Preliminary results from market maker simulations demonstrate key findings: (1) sentiment-spread correlation of 0.55, confirming market makers widen quotes during sentiment extremes; (2) uncertainty-spread correlation of 0.72, validating that epistemic uncertainty drives spread adjustment; and (3) proper sentiment differentiation across bullish (+0.79), neutral (+0.23), and bearish ($-0.92$) texts with appropriate uncertainty decomposition.

The framework enables investigation of how sentiment-driven trading affects market microstructure properties including bid-ask spreads, volatility clustering, and flash crash dynamics.

\vspace{0.5em}
\noindent\textbf{Keywords:} agent-based modeling, market microstructure, sentiment analysis, Monte Carlo dropout, cryptocurrency, Kafka, uncertainty quantification

\vspace{0.3em}
\noindent\textbf{JEL Codes:} G12 (Asset Pricing), G14 (Information and Market Efficiency), C63 (Computational Techniques; Simulation Modeling)
\end{abstract}

\vspace{1em}
\metadatabox{\paperdoi}

\clearpage
\twocolumn
\setstretch{1.15}

\section{Introduction}
\label{sec:introduction}

Cryptocurrency markets exhibit distinctive microstructure properties that differ from traditional equity markets: 24/7 continuous trading, extreme volatility clustering, flash crashes, and apparent sensitivity to social media sentiment \citep{bouri2017hedge}. Understanding how these properties emerge from trader behavior requires models that capture both the heterogeneity of market participants and the information channels they respond to.

This paper presents a methodological framework for agent-based modeling (ABM) of cryptocurrency microstructure with integrated real-time sentiment analysis. The key innovation is treating sentiment as an \textit{uncertain signal} rather than a deterministic input: our Monte Carlo Dropout sentiment analyzer provides not just sentiment scores but quantified uncertainty, enabling heterogeneous agent responses to sentiment information.

\subsection{Motivation}

Traditional market microstructure models assume rational agents with well-defined information structures \citep{glosten1985bid, kyle1985continuous}. Cryptocurrency markets challenge this assumption:

\begin{itemize}[leftmargin=1.5em, topsep=0pt, itemsep=2pt]
    \item \textbf{Sentiment dominance:} Price movements correlate with social media activity \citep{abraham2018cryptocurrency}
    \item \textbf{Uncertainty about uncertainty:} Traders disagree about the informativeness of sentiment signals
    \item \textbf{Regime dependence:} Bull and bear markets exhibit different microstructure properties
\end{itemize}

Our framework addresses these challenges by:

\begin{enumerate}[leftmargin=1.5em, topsep=0pt, itemsep=2pt]
    \item Quantifying sentiment uncertainty using Monte Carlo Dropout
    \item Specifying heterogeneous agent responses to uncertain sentiment
    \item Detecting market regimes via Dynamic Factor Models
    \item Enabling real-time simulation with streaming data infrastructure
\end{enumerate}

\subsection{Contributions}

This methodology paper contributes:

\begin{enumerate}[leftmargin=1.5em, topsep=0pt, itemsep=2pt]
    \item An uncertainty-aware sentiment analysis pipeline using Monte Carlo Dropout
    \item Agent specifications for market makers, informed traders, noise traders, and arbitrageurs with sentiment-responsive behavior
    \item A distributed computing architecture for real-time ABM simulation
    \item Calibration methodology for matching simulated to empirical stylized facts
\end{enumerate}

\section{Related Work}
\label{sec:related}

\subsection{Agent-Based Market Models}

Agent-based computational economics has a rich history of market simulation \citep{lebaron2006agent}. The Santa Fe Artificial Stock Market \citep{palmer1994artificial} demonstrated emergence of realistic market dynamics from simple agent rules. Subsequent work has explored order book dynamics \citep{cont2010stochastic}, flash crashes \citep{paddrik2012agent}, and market design \citep{farmer2005predictive}.

\subsection{Cryptocurrency Market Microstructure}

\citet{makarov2020trading} document significant price dislocations across crypto exchanges, suggesting fragmented liquidity. \citet{hautsch2018limits} analyze Bitcoin order flow dynamics. However, existing work largely treats sentiment as exogenous; we integrate sentiment as an endogenous information channel.

\subsection{Uncertainty-Aware Sentiment Analysis}

Standard sentiment classifiers provide point estimates without uncertainty quantification. Monte Carlo Dropout \citep{gal2016dropout} enables approximate Bayesian inference by running multiple forward passes with dropout enabled at inference time, producing a distribution of predictions that captures model uncertainty.

\section{Framework Architecture}
\label{sec:architecture}

\subsection{System Overview}

The framework consists of four layers:

\begin{enumerate}[leftmargin=1.5em, topsep=0pt, itemsep=2pt]
    \item \textbf{Data Ingestion:} Real-time feeds from Binance (orderbook) and Reddit (social)
    \item \textbf{Feature Engineering:} Sentiment scoring with uncertainty, microstructure metrics
    \item \textbf{Regime Detection:} Dynamic Factor Model for market state
    \item \textbf{Simulation:} Mesa-based ABM with heterogeneous agents
\end{enumerate}

\subsection{Data Ingestion Layer}

\subsubsection{Binance WebSocket Client}

Orderbook depth streams provide 100ms updates for:

\begin{itemize}[leftmargin=1.5em, topsep=0pt, itemsep=2pt]
    \item Best bid/offer prices and sizes
    \item Order book imbalance: $\text{imb} = \frac{V_{bid} - V_{ask}}{V_{bid} + V_{ask}}$
    \item Mid-price: $p_{mid} = \frac{p_{bid} + p_{ask}}{2}$
    \item Spread: $s = \frac{p_{ask} - p_{bid}}{p_{mid}}$
\end{itemize}

\subsubsection{Reddit API Client}

Streaming ingestion from 7 cryptocurrency subreddits:
\begin{itemize}[leftmargin=1.5em, topsep=0pt, itemsep=2pt]
    \item r/CryptoCurrency, r/Bitcoin, r/ethereum
    \item r/CryptoMarkets, r/altcoin
    \item r/binance, r/defi
\end{itemize}

Posts and comments are published to Kafka topic for downstream processing.

\subsection{Sentiment Analysis with Uncertainty}

\subsubsection{Model Architecture}

We use CryptoBERT \citep{cryptobert2024}, a RoBERTa-based model fine-tuned on 3.2 million cryptocurrency social media posts from StockTwits, with Monte Carlo Dropout for uncertainty quantification. The model was trained with balanced labels: Bearish (0), Neutral (1), Bullish (2).

\textbf{Monte Carlo Dropout procedure:}

\begin{enumerate}[leftmargin=1.5em, topsep=0pt, itemsep=2pt]
    \item Enable dropout at inference time
    \item Run $T=50$ forward passes
    \item Compute mean prediction: $\bar{y} = \frac{1}{T}\sum_{t=1}^T f_\theta(x)_t$
    \item Compute epistemic uncertainty: $\sigma_{epi}^2 = \frac{1}{T}\sum_{t=1}^T (f_\theta(x)_t - \bar{y})^2$
\end{enumerate}

\subsubsection{Uncertainty Decomposition}

Total uncertainty decomposes into:

\begin{equation}
\sigma_{total}^2 = \underbrace{\sigma_{epi}^2}_{\text{model uncertainty}} + \underbrace{H(\bar{y})}_{\text{aleatoric uncertainty}}
\end{equation}

where $H(\bar{y})$ is the Shannon entropy of the mean prediction, capturing irreducible uncertainty from ambiguous text.

\subsubsection{EWMA Smoothing}

Raw sentiment is smoothed using exponentially weighted moving average:

\begin{equation}
s_t = \alpha \cdot s_{raw,t} + (1-\alpha) \cdot s_{t-1}
\end{equation}

with $\alpha = 0.1$ corresponding to approximately 5-minute half-life.

\subsection{Agent Specifications}

We implement four agent types using the Mesa framework.

\subsubsection{Market Makers}

Market makers provide liquidity by quoting bid and ask prices:

\begin{align}
p_{bid} &= p_{mid} - \frac{s}{2} - \gamma \cdot Q - \delta \cdot \sigma_{epi} \\
p_{ask} &= p_{mid} + \frac{s}{2} + \gamma \cdot Q + \delta \cdot \sigma_{epi}
\end{align}

where $Q$ is inventory, $\gamma$ is inventory aversion, and $\delta$ scales spread widening with sentiment uncertainty.

\textbf{Key behavior:} Market makers widen spreads when sentiment uncertainty is high, reflecting increased adverse selection risk.

\subsubsection{Informed Traders}

Informed traders trade on sentiment signals when confidence is high:

\begin{equation}
\text{trade} =
\begin{cases}
\text{buy } V & \text{if } s_t > \tau \text{ and } \sigma_{epi} < \bar{\sigma} \\
\text{sell } V & \text{if } s_t < -\tau \text{ and } \sigma_{epi} < \bar{\sigma} \\
\text{hold} & \text{otherwise}
\end{cases}
\end{equation}

where $\tau$ is sentiment threshold and $\bar{\sigma}$ is maximum acceptable uncertainty.

\textbf{Key behavior:} Informed traders only act when sentiment signal is strong AND model is confident.

\subsubsection{Noise Traders}

Noise traders arrive according to a Poisson process with intensity $\lambda$ and submit market orders:

\begin{equation}
\text{direction} \sim \text{Bernoulli}(0.5 + \beta \cdot s_t)
\end{equation}

where $\beta$ controls sentiment influence on noise trader direction.

\textbf{Key behavior:} Noise traders are weakly influenced by sentiment regardless of uncertainty.

\subsubsection{Arbitrageurs}

Arbitrageurs exploit price dislocations:

\begin{equation}
\text{trade} =
\begin{cases}
\text{buy} & \text{if } p < p_{fair} - \epsilon \\
\text{sell} & \text{if } p > p_{fair} + \epsilon \\
\text{hold} & \text{otherwise}
\end{cases}
\end{equation}

where $p_{fair}$ is estimated from cross-exchange prices or fundamental indicators.

\textbf{Key behavior:} Arbitrageurs are sentiment-agnostic, responding only to price dislocations.

\subsection{Regime Detection}

A Dynamic Factor Model extracts latent market state:

\begin{equation}
\mathbf{y}_t = \mathbf{\Lambda} \mathbf{f}_t + \mathbf{e}_t
\end{equation}

where $\mathbf{y}_t$ is observed microstructure metrics, $\mathbf{f}_t$ is latent factors, and $\mathbf{\Lambda}$ is factor loadings.

Factor values are used to switch agent parameters between bull/bear regimes, enabling adaptive behavior.

\section{Infrastructure}
\label{sec:infrastructure}

\subsection{Distributed Architecture}

The framework deploys on Kubernetes (K3s) with:

\begin{itemize}[leftmargin=1.5em, topsep=0pt, itemsep=2pt]
    \item \textbf{Kafka:} Event streaming for data ingestion
    \item \textbf{TimescaleDB:} Time-series storage with continuous aggregates
    \item \textbf{GPU nodes:} Sentiment model inference
    \item \textbf{Mesa simulation:} Parallelized agent execution
\end{itemize}

\subsection{Data Flow}

\begin{enumerate}[leftmargin=1.5em, topsep=0pt, itemsep=2pt]
    \item Binance WebSocket $\rightarrow$ Kafka topic ``order-books''
    \item Reddit API $\rightarrow$ Kafka topic ``reddit-posts''
    \item Sentiment service consumes reddit-posts, produces ``sentiment-ticks''
    \item Simulation service consumes all topics, runs ABM, writes to TimescaleDB
\end{enumerate}

\subsection{Performance Targets}

\begin{itemize}[leftmargin=1.5em, topsep=0pt, itemsep=2pt]
    \item Order book updates: 10/second
    \item Sentiment processing: <100ms latency
    \item Kafka end-to-end lag: <500ms
    \item Simulation step: <500ms for 1000 agents
\end{itemize}

\section{Calibration Methodology}
\label{sec:calibration}

\subsection{Stylized Facts}

The simulation should reproduce empirical stylized facts \citep{cont2001empirical}:

\begin{enumerate}[leftmargin=1.5em, topsep=0pt, itemsep=2pt]
    \item \textbf{Volatility clustering:} $\text{ACF}(|r|, \text{lag}=10) > 0.1$
    \item \textbf{Fat tails:} Kurtosis $> 3$
    \item \textbf{Spread mean-reversion:} ADF test $p < 0.05$
    \item \textbf{Volume-volatility correlation:} $\rho(V, |r|) > 0$
\end{enumerate}

\subsection{Parameter Calibration}

Agent parameters are calibrated via:

\begin{enumerate}[leftmargin=1.5em, topsep=0pt, itemsep=2pt]
    \item Collect 1 week empirical Binance data
    \item Compute target stylized facts
    \item Grid search over agent parameters
    \item Select parameters minimizing stylized fact deviation
\end{enumerate}

\subsection{Validation}

\begin{itemize}[leftmargin=1.5em, topsep=0pt, itemsep=2pt]
    \item K-S test for return distribution matching
    \item Cross-correlation comparison: simulated vs.\ empirical
    \item Out-of-sample prediction: simulate next day, compare
\end{itemize}

\section{Research Questions}
\label{sec:questions}

The framework enables investigation of:

\begin{enumerate}[leftmargin=1.5em, topsep=0pt, itemsep=2pt]
    \item How does sentiment uncertainty affect market maker spread-setting?
    \item Do informed traders provide price discovery or amplify noise?
    \item What agent composition produces realistic flash crash dynamics?
    \item How do sentiment shocks propagate through market microstructure?
\end{enumerate}

\section{Implementation Status}
\label{sec:status}

\textbf{Completed:}
\begin{itemize}[leftmargin=1.5em, topsep=0pt, itemsep=2pt]
    \item Data ingestion layer (Binance, Reddit clients)
    \item Monte Carlo Dropout sentiment analyzer
    \item Infrastructure configuration (Docker, Kafka, TimescaleDB)
\end{itemize}

\textbf{In Progress:}
\begin{itemize}[leftmargin=1.5em, topsep=0pt, itemsep=2pt]
    \item Agent implementations (market maker, informed, noise, arbitrageur)
    \item Order book matching engine
    \item Mesa simulation environment
\end{itemize}

\textbf{Planned:}
\begin{itemize}[leftmargin=1.5em, topsep=0pt, itemsep=2pt]
    \item DFM regime detection
    \item Calibration pipeline
    \item Dashboard visualization
\end{itemize}

\section{Preliminary Results}
\label{sec:results}

We present preliminary results from sentiment analysis and market maker simulation components to validate core framework functionality.

\subsection{Sentiment Analysis Validation}

We evaluate CryptoBERT with MC Dropout on 12 representative cryptocurrency texts spanning bullish, neutral, and bearish sentiment. Table~\ref{tab:sentiment} presents results.

\begin{table}[h]
\centering
\caption{Sentiment Analysis Results with Uncertainty Quantification}
\label{tab:sentiment}
\small
\begin{tabular}{p{3.8cm}ccc}
\toprule
\textbf{Text (truncated)} & \textbf{Sent.} & \textbf{Epist.} & \textbf{Aleat.} \\
\midrule
BTC ETF approved! Huge... & +0.65 & 0.002 & 0.68 \\
Just bought more BTC... & +0.79 & 0.001 & 0.51 \\
Sideways around 43k... & +0.23 & 0.004 & 0.54 \\
SEC meeting next week... & $-$0.36 & 0.059 & 0.88 \\
FTX collapse, crypto scam... & $-$0.92 & 0.000 & 0.29 \\
SEC suing, crypto winter... & $-$0.90 & 0.011 & 0.34 \\
Bull trap? Staying cautious... & +0.10 & 0.001 & 0.36 \\
\bottomrule
\end{tabular}
\end{table}

Key observations:
\begin{itemize}[leftmargin=1.5em, topsep=0pt, itemsep=2pt]
    \item \textbf{Sentiment differentiation:} Model correctly identifies bullish (+0.65 to +0.79), neutral (+0.10 to +0.23), and bearish ($-$0.36 to $-$0.92) texts
    \item \textbf{Epistemic uncertainty:} Very low (0.001--0.011) for most texts, indicating model confidence; higher (0.059) for regulatory news suggesting domain uncertainty
    \item \textbf{Aleatoric uncertainty:} Higher for ambiguous texts (0.88 for SEC meeting) than clear sentiment (0.29 for FTX collapse)
\end{itemize}

\subsection{Market Maker Simulation}

We simulate 2,000 timesteps of market maker behavior responding to synthetic sentiment series with regime changes. Figure~\ref{fig:results} presents four-panel visualization of the simulation dynamics.

\begin{figure}[h]
\centering
\includegraphics[width=\columnwidth]{figures/preliminary_results.pdf}
\caption{Preliminary simulation results: (A) Sentiment scores with uncertainty bars; (B) Epistemic vs.\ aleatoric uncertainty decomposition; (C) Market maker spread response to sentiment; (D) Inventory dynamics.}
\label{fig:results}
\end{figure}

Table~\ref{tab:summary_stats} presents summary statistics for the simulation.

\begin{table}[h]
\centering
\caption{Simulation Summary Statistics ($n = 2,000$ timesteps)}
\label{tab:summary_stats}
\small
\begin{tabular}{lrrr}
\toprule
\textbf{Statistic} & \textbf{Log Return} & \textbf{Spread (bps)} & \textbf{Sentiment} \\
\midrule
Mean & 0.000066 & 0.1085 & 0.0472 \\
Std Dev & 0.000714 & 0.0068 & 0.3932 \\
Min & $-$0.00278 & 0.0914 & $-$1.0000 \\
Max & 0.00233 & 0.1308 & 1.0000 \\
Skewness & $-$0.124 & $-$0.161 & $-$0.292 \\
Kurtosis & 0.016 & $-$0.352 & $-$0.294 \\
\bottomrule
\end{tabular}
\end{table}

\subsection{Return Distribution Analysis}

We examine the distributional properties of simulated returns to assess alignment with empirical stylized facts. Figure~\ref{fig:return_dist} presents the return distribution with Q-Q plot against the normal distribution.

\begin{figure}[h]
\centering
\includegraphics[width=\columnwidth]{figures/return_distribution.pdf}
\caption{Return distribution analysis: (A) Histogram of log returns with fitted normal distribution; (B) Q-Q plot against normal distribution showing slight tail deviation.}
\label{fig:return_dist}
\end{figure}

The Jarque-Bera test statistic of 5.15 ($p = 0.076$) marginally fails to reject normality at the 5\% level. Excess kurtosis of 0.016 indicates absence of pronounced fat tails in this baseline simulation---a limitation reflecting the single market maker configuration without informed trader interactions. Negative skewness ($-$0.124) suggests slight asymmetry toward negative returns.

\subsection{Time-Series Diagnostics}

We conduct formal stationarity and autocorrelation tests to validate time-series properties. Table~\ref{tab:diagnostics} summarizes results.

\begin{table}[h]
\centering
\caption{Time-Series Diagnostic Tests}
\label{tab:diagnostics}
\small
\begin{tabular}{lccc}
\toprule
\textbf{Test} & \textbf{Statistic} & \textbf{p-value} & \textbf{Conclusion} \\
\midrule
ADF (Spread) & $-$2.62 & 0.089 & Non-stationary \\
KPSS (Spread) & 1.00 & $<$0.01 & Non-stationary \\
Jarque-Bera (Returns) & 5.15 & 0.076 & Approx.\ Normal \\
\bottomrule
\end{tabular}
\end{table}

The spread series exhibits non-stationarity by both ADF ($p = 0.089$) and KPSS ($p < 0.01$) tests, indicating potential regime-dependent behavior. Figure~\ref{fig:acf} presents autocorrelation analysis.

\begin{figure}[h]
\centering
\includegraphics[width=\columnwidth]{figures/acf_comparison.pdf}
\caption{Autocorrelation analysis: (A) ACF of returns showing significant positive autocorrelation from momentum effects; (B) ACF of absolute returns. Both panels include 95\% confidence intervals.}
\label{fig:acf}
\end{figure}

Ljung-Box tests indicate significant autocorrelation in both returns ($Q_{10} = 1467.3$, $p < 0.001$) and absolute returns ($Q_{10} = 37.4$, $p < 0.001$). The strong autocorrelation in returns reflects the momentum dynamics in the market maker's quote adjustment process.

\subsection{Regime-Conditional Dynamics}

We classify each timestep into bullish ($s_t > 0.2$), neutral ($|s_t| \leq 0.2$), or bearish ($s_t < -0.2$) regimes based on sentiment scores. Table~\ref{tab:regime_stats} presents regime-conditional statistics.

\begin{table}[h]
\centering
\caption{Regime-Conditional Statistics}
\label{tab:regime_stats}
\small
\begin{tabular}{lrrrrr}
\toprule
\textbf{Regime} & \textbf{N} & \textbf{\% Time} & \textbf{$\bar{s}$ (bps)} & \textbf{$\sigma_r$} & \textbf{$\bar{\sigma}_u$} \\
\midrule
Bullish & 488 & 24.4\% & 10.57 & 0.063\% & 0.247 \\
Neutral & 1,144 & 57.2\% & 11.19 & 0.065\% & 0.301 \\
Bearish & 368 & 18.4\% & 10.17 & 0.065\% & 0.245 \\
\bottomrule
\multicolumn{6}{l}{\footnotesize $\bar{s}$: mean spread; $\sigma_r$: return volatility; $\bar{\sigma}_u$: mean total uncertainty.}
\end{tabular}
\end{table}

Interestingly, spreads are \textit{widest} during neutral regimes (11.19 bps) and narrowest during bearish regimes (10.17 bps). This counter-intuitive result reflects the higher uncertainty during neutral regimes ($\bar{\sigma}_u = 0.301$) compared to extreme sentiment regimes.

Figure~\ref{fig:regime_dynamics} visualizes the regime dynamics and duration distribution.

\begin{figure}[h]
\centering
\includegraphics[width=\columnwidth]{figures/regime_dynamics.pdf}
\caption{Regime dynamics: (A) Sentiment time series with regime identification (green: bullish, gray: neutral, red: bearish); (B) Spread response to regimes; (C) Distribution of regime durations by type.}
\label{fig:regime_dynamics}
\end{figure}

Table~\ref{tab:transitions} presents the regime transition probability matrix.

\begin{table}[h]
\centering
\caption{Regime Transition Matrix $P(\text{regime}_{t+1} | \text{regime}_t)$}
\label{tab:transitions}
\small
\begin{tabular}{lccc}
\toprule
& \textbf{$\to$ Bullish} & \textbf{$\to$ Neutral} & \textbf{$\to$ Bearish} \\
\midrule
\textbf{Bullish $\to$} & 0.992 & 0.004 & 0.004 \\
\textbf{Neutral $\to$} & 0.003 & 0.995 & 0.002 \\
\textbf{Bearish $\to$} & 0.003 & 0.008 & 0.989 \\
\bottomrule
\end{tabular}
\end{table}

High diagonal persistence ($>$98.9\%) indicates regime stability, with regimes lasting 92--191 steps on average. Bearish regimes have shortest mean duration (92 steps) compared to neutral (191 steps).

\subsection{Uncertainty Decomposition}

We analyze the decomposition of total uncertainty into epistemic and aleatoric components. Table~\ref{tab:correlation} presents the correlation structure.

\begin{table}[h]
\centering
\caption{Correlation Matrix (Key Variables)}
\label{tab:correlation}
\small
\begin{tabular}{lrrrr}
\toprule
& \textbf{Sent.} & \textbf{$\sigma_e$} & \textbf{$\sigma_a$} & \textbf{Spread} \\
\midrule
Sentiment & 1.000 & $-$0.013 & $-$0.008 & 0.085 \\
$\sigma_{\text{epistemic}}$ & $-$0.013 & 1.000 & 0.559 & \textbf{0.496} \\
$\sigma_{\text{aleatoric}}$ & $-$0.008 & 0.559 & 1.000 & \textbf{0.612} \\
Spread & 0.085 & \textbf{0.496} & \textbf{0.612} & 1.000 \\
\bottomrule
\end{tabular}
\end{table}

Key findings:
\begin{itemize}[leftmargin=1.5em, topsep=0pt, itemsep=2pt]
    \item \textbf{Aleatoric dominance:} Aleatoric uncertainty comprises 81.6\% of total uncertainty (mean: 0.227 vs.\ 0.051 epistemic)
    \item \textbf{Spread correlations:} Aleatoric ($\rho = 0.612$) stronger predictor than epistemic ($\rho = 0.496$)
    \item \textbf{Sentiment-spread:} Direct sentiment-spread correlation weak ($\rho = 0.085$) compared to uncertainty effects
\end{itemize}

Figure~\ref{fig:uncertainty} visualizes the uncertainty decomposition.

\begin{figure}[h]
\centering
\includegraphics[width=\columnwidth]{figures/uncertainty_decomposition.pdf}
\caption{Uncertainty decomposition: (A) Stacked area plot of epistemic (blue) and aleatoric (red) uncertainty over time; (B) Scatter plot showing spread response to each uncertainty component with linear fits.}
\label{fig:uncertainty}
\end{figure}

The stronger aleatoric-spread correlation suggests market makers respond more to inherent signal ambiguity than to model uncertainty. This has practical implications: sentiment signals with high aleatoric uncertainty warrant wider spreads regardless of model confidence.

\section{Discussion}
\label{sec:discussion}

\subsection{Validation of Core Hypotheses}

The expanded preliminary results provide nuanced support for our theoretical framework:

\textbf{H1: Uncertainty drives microstructure adjustment.} Counter-intuitively, direct sentiment-spread correlation is weak ($\rho = 0.085$), while uncertainty-spread correlation is strong (total: $\rho = 0.637$; aleatoric: $\rho = 0.612$). This suggests market makers respond primarily to \textit{signal quality} rather than signal direction---consistent with adverse selection theory where informed trading risk matters more than market direction.

\textbf{H2: Aleatoric uncertainty dominates.} Aleatoric (inherent) uncertainty comprises 81.6\% of total uncertainty and shows stronger spread correlation than epistemic (model) uncertainty. This validates our decomposition approach: market makers should weight inherent signal ambiguity over model confidence when adjusting quotes.

\textbf{H3: Regime dynamics reveal counter-intuitive patterns.} Spreads are \textit{widest} during neutral regimes (11.19 bps) despite sentiment extremes during bullish/bearish periods. This reflects higher uncertainty during neutral regimes ($\bar{\sigma}_u = 0.301$ vs.\ 0.247 in bullish). The finding supports uncertainty-weighted rather than sentiment-weighted market making.

\subsection{Comparison to Traditional Models}

Our framework extends traditional market microstructure models \citep{glosten1985bid, kyle1985continuous} by:
\begin{enumerate}[leftmargin=1.5em, topsep=0pt, itemsep=2pt]
    \item Incorporating real-time social sentiment as an information channel
    \item Quantifying uncertainty in the sentiment signal
    \item Enabling heterogeneous agent responses to uncertain information
\end{enumerate}

The Avellaneda-Stoikov \citeyearpar{avellaneda2008high} market-making model provides the foundation, extended with sentiment-dependent spread adjustment and uncertainty premium.

\subsection{Limitations}

Several limitations apply to preliminary results:

\begin{enumerate}[leftmargin=1.5em, topsep=0pt, itemsep=2pt]
    \item \textbf{Synthetic dynamics:} Current simulation uses synthetic sentiment series rather than real Reddit data
    \item \textbf{Single agent type:} Only market maker agents implemented; informed traders, noise traders, and arbitrageurs pending
    \item \textbf{No order book matching:} Simplified quote generation without full limit order book dynamics
    \item \textbf{Absence of fat tails:} Return kurtosis (0.016) approximates normal distribution; real markets exhibit excess kurtosis $>$ 3. This limitation reflects the single-agent configuration without informed trader interactions
    \item \textbf{Strong autocorrelation:} Returns exhibit significant positive autocorrelation reflecting momentum dynamics, rather than empirical near-zero autocorrelation
    \item \textbf{Non-stationary spreads:} Spread series fails ADF/KPSS stationarity tests, suggesting regime-dependent dynamics requiring further investigation
\end{enumerate}

\subsection{Implications for Cryptocurrency Markets}

The strong uncertainty-spread correlation ($\rho = 0.637$) and aleatoric dominance have practical implications:
\begin{itemize}[leftmargin=1.5em, topsep=0pt, itemsep=2pt]
    \item Market makers should monitor sentiment \textit{uncertainty}, not just sentiment level---the direct sentiment-spread correlation is only 0.085
    \item Neutral sentiment periods with high uncertainty warrant \textit{wider} spreads than extreme sentiment periods with clear signals
    \item Aleatoric uncertainty (81.6\% of total) provides more actionable information than epistemic uncertainty for spread adjustment
    \item Regime persistence ($>$98.9\%) suggests stable trading regimes once established, reducing regime switching costs
\end{itemize}

\subsection{Next Steps}

Completing the full simulation requires:
\begin{enumerate}[leftmargin=1.5em, topsep=0pt, itemsep=2pt]
    \item Implementing remaining agent types (informed, noise, arbitrageur)
    \item Building Mesa-based order book matching engine
    \item Collecting 1-week empirical Binance + Reddit data
    \item Calibrating agent parameters to match stylized facts
    \item Running shock scenarios (sentiment crashes, flash crash emergence)
\end{enumerate}

\section{Conclusion}
\label{sec:conclusion}

This paper presents a methodological framework for agent-based modeling of cryptocurrency market microstructure with uncertainty-aware sentiment integration. The key innovation is treating sentiment as a noisy signal with decomposed uncertainty (epistemic vs.\ aleatoric), enabling nuanced modeling of market maker responses to social media information quality.

Preliminary results from 2,000 timesteps reveal several findings:
\begin{itemize}[leftmargin=1.5em, topsep=0pt, itemsep=2pt]
    \item \textbf{Uncertainty dominates sentiment:} Direct sentiment-spread correlation is weak ($\rho = 0.085$) while total uncertainty-spread correlation is strong ($\rho = 0.637$)
    \item \textbf{Aleatoric uncertainty is key:} Aleatoric (inherent) uncertainty comprises 81.6\% of total and shows stronger spread correlation ($\rho = 0.612$) than epistemic uncertainty ($\rho = 0.496$)
    \item \textbf{Counter-intuitive regime effects:} Spreads are widest during neutral sentiment regimes (11.19 bps) due to higher uncertainty, not during sentiment extremes
    \item \textbf{High regime persistence:} Transition probabilities exceed 98.9\% on diagonal, indicating stable trading regimes
\end{itemize}

The framework combines modern NLP (Monte Carlo Dropout for uncertainty), domain-specific models (CryptoBERT for cryptocurrency sentiment), and agent-based simulation into an integrated pipeline. Current limitations include absence of fat tails (kurtosis = 0.016) and significant return autocorrelation---both reflecting the single market maker configuration without informed trader interactions.

Future research will implement additional agent types, calibrate parameters to empirical stylized facts, and investigate how sentiment shocks propagate through multi-agent market microstructure.

\clearpage
\onecolumn

\begin{thebibliography}{99}

\bibitem[Avellaneda \& Stoikov(2008)]{avellaneda2008high}
Avellaneda, M., \& Stoikov, S. (2008).
\newblock High-frequency trading in a limit order book.
\newblock \textit{Quantitative Finance}, 8(3), 217--224.

\bibitem[Abraham et al.(2018)]{abraham2018cryptocurrency}
Abraham, J., Higdon, D., Nelson, J., \& Ibarra, J. (2018).
\newblock Cryptocurrency price prediction using tweet volumes and sentiment analysis.
\newblock \textit{SMU Data Science Review}, 1(3), 1.

\bibitem[Bouri et al.(2017)]{bouri2017hedge}
Bouri, E., Molnar, P., Azzi, G., Roubaud, D., \& Hagfors, L. I. (2017).
\newblock On the hedge and safe haven properties of Bitcoin.
\newblock \textit{Research in International Business and Finance}, 41, 88--98.

\bibitem[Cont(2001)]{cont2001empirical}
Cont, R. (2001).
\newblock Empirical properties of asset returns: Stylized facts and statistical issues.
\newblock \textit{Quantitative Finance}, 1(2), 223--236.

\bibitem[Cont et al.(2010)]{cont2010stochastic}
Cont, R., Stoikov, S., \& Talreja, R. (2010).
\newblock A stochastic model for order book dynamics.
\newblock \textit{Operations Research}, 58(3), 549--563.

\bibitem[ElKulako(2024)]{cryptobert2024}
ElKulako (2024).
\newblock CryptoBERT: A language model for cryptocurrency sentiment.
\newblock Hugging Face. \url{https://huggingface.co/ElKulako/cryptobert}

\bibitem[Farmer \& Foley(2009)]{farmer2005predictive}
Farmer, J. D., \& Foley, D. (2009).
\newblock The economy needs agent-based modelling.
\newblock \textit{Nature}, 460(7256), 685--686.

\bibitem[Gal \& Ghahramani(2016)]{gal2016dropout}
Gal, Y., \& Ghahramani, Z. (2016).
\newblock Dropout as a Bayesian approximation: Representing model uncertainty in deep learning.
\newblock \textit{ICML}, 1050--1059.

\bibitem[Glosten \& Milgrom(1985)]{glosten1985bid}
Glosten, L. R., \& Milgrom, P. R. (1985).
\newblock Bid, ask and transaction prices in a specialist market.
\newblock \textit{Journal of Financial Economics}, 14(1), 71--100.

\bibitem[Hautsch et al.(2018)]{hautsch2018limits}
Hautsch, N., Scheuch, C., \& Voigt, S. (2018).
\newblock Limits to arbitrage in markets with stochastic settlement latency.
\newblock \textit{Working Paper}.

\bibitem[Kyle(1985)]{kyle1985continuous}
Kyle, A. S. (1985).
\newblock Continuous auctions and insider trading.
\newblock \textit{Econometrica}, 53(6), 1315--1335.

\bibitem[LeBaron(2006)]{lebaron2006agent}
LeBaron, B. (2006).
\newblock Agent-based computational finance.
\newblock \textit{Handbook of Computational Economics}, 2, 1187--1233.

\bibitem[Makarov \& Schoar(2020)]{makarov2020trading}
Makarov, I., \& Schoar, A. (2020).
\newblock Trading and arbitrage in cryptocurrency markets.
\newblock \textit{Journal of Financial Economics}, 135(2), 293--319.

\bibitem[Paddrik et al.(2012)]{paddrik2012agent}
Paddrik, M., Hayes, R., Todd, A., Yang, S., Beling, P., \& Scherer, W. (2012).
\newblock An agent based model of the E-Mini S\&P 500 applied to flash crash analysis.
\newblock \textit{IEEE Conference on Computational Intelligence for Financial Engineering}.

\bibitem[Palmer et al.(1994)]{palmer1994artificial}
Palmer, R. G., Arthur, W. B., Holland, J. H., LeBaron, B., \& Tayler, P. (1994).
\newblock Artificial economic life: A simple model of a stockmarket.
\newblock \textit{Physica D}, 75(1-3), 264--274.

\end{thebibliography}

\end{document}
